% Generated draft: requires mathematical and source review.
\documentclass[11pt]{article}
\usepackage[margin=25mm]{geometry}
\usepackage{fontspec}
\setmainfont{Latin Modern Roman}
\usepackage{amsmath,amssymb,mathtools,unicode-math}
\setmathfont{Latin Modern Math}
\usepackage{enumitem,needspace,xurl,hyperref}
\setlength{\parindent}{0pt}
\setlength{\parskip}{5pt}
\setlength{\emergencystretch}{3em}
\newcommand{\N}{\mathbb{N}}
\newcommand{\Z}{\mathbb{Z}}
\newcommand{\Q}{\mathbb{Q}}
\newcommand{\R}{\mathbb{R}}
\newcommand{\C}{\mathbb{C}}
\newcommand{\F}{\mathbb{F}}
\newcommand{\A}{\mathbb{A}}
\newcommand{\PP}{\mathbb{P}}
\newcommand{\E}{\mathbb{E}}
\newcommand{\eps}{\varepsilon}
\newcommand{\dd}{\,\mathrm{d}}
\begin{document}
\section*{UnsolvedMath ID: 30005972}
\textbf{Source problem code:} catalogue entry
\subsection{Definitions and mathematical statement}
Let \(K\) be a discretely valued field with valuation \(v\), valuation ring \(R_K\), and residue field \(k_K\), and let \(Y/K\) be the smooth projective curve in the supplied record. For a finite extension \(L/K\), choose an extension \(w\) of \(v\), write \(R_L\) for its valuation ring and \(k_L\) for its residue field, and set
\[
Y_L=Y\times_{\operatorname{Spec}K}\operatorname{Spec}L.
\]
A semistable model means a proper, flat, finitely presented \(R_L\)-scheme \(\mathcal Y\) of relative dimension one, together with an identification \(\mathcal Y_L\cong Y_L\), whose geometric special fiber
\[
\mathcal Y_{\overline s}=\bigl(\mathcal Y\times_{R_L}k_L\bigr)\times_{k_L}\overline{k_L}
\]
is reduced and has only ordinary double points as singularities. Write \(\mathcal Y_s=\mathcal Y\times_{R_L}k_L\).

The requested construction is an explicit method, valid in arbitrary genus, that takes \(Y/K\) as input and computes a finite extension \(L/K\), a semistable model \(\mathcal Y/R_L\) of \(Y_L\), and the action described in the source as the action of \(\operatorname{Gal}(L/K)\) on \(\mathcal Y_s\). If \(Y\) is geometrically connected, its genus is \(g=\dim_K H^1(Y,\mathcal O_Y)\). The source describes the construction as trivial for \(g=0\), well known for \(g=1\), and only partially available in special cases for \(g\geq2\); these are supplied qualifications, not independently verified conclusions.

The Galois-action requirement needs further specification. The record requests only a finite extension, without explicitly requiring it to be Galois. Even for a finite Galois extension, an action on the chosen special fiber requires an invariant extension of the valuation and a compatible extension of the generic-fiber descent action to suitable model data. Such an induced action is generally semilinear: an element \(\sigma\) acts on the special fiber over its induced automorphism \(\overline\sigma\) of \(k_L\), rather than necessarily as a \(k_L\)-linear automorphism. The source does not specify whether the model itself must carry this action, how the action is to be represented, or whether a geometric special fiber is intended. These conditions are unresolved parts of the requested formulation, not additional hypotheses silently imposed here.

The field assumptions and computational input are also unresolved. The supplied record does not specify completeness or henselianity of \(K\), conditions on its residue field or characteristic, geometric connectedness of \(Y\), or an effective presentation of the field and curve. Consequently this is a statement of the intended construction problem with explicit qualifications, rather than a fully specified algorithmic assertion for every abstract discretely valued field.
\subsection{Short English statement}
Develop an explicit method to compute semistable reduction of a smooth projective curve after a finite field extension, together with the corresponding Galois action on the special fiber, in arbitrary genus. The supplied record reports partial solutions, but leaves the field hypotheses, effective input, and precise Galois-action requirement insufficiently specified.
\subsection{Sources}
\begin{itemize}
\item Local UnsolvedMath registry, supplied complete record: permanent ID \(30005972\), source problem code \texttt{catalogue entry}, title \emph{Explicit Semistable Reduction and Galois Actions for Curves}; collection: Oberwolfach Reports Open Problems.
\item Supplied report citation: \emph{Explicit Methods in Number Theory (2025)}, Oberwolfach Reports, DOI \texttt{10.4171/owr/2024/40}, pp.\ 2309--2368. \url{https://doi.org/10.4171/owr/2024/40}. The excerpt identifies Problem 1 and the contribution \emph{Semistable reduction of covers of curves of degree \(p\)}, Stefan Wewers, joint work with Ole Ossen. The registry supplies no separate problem proposer and records proposed year 2024.
\item Supplied literature source: Dokchitser, \emph{A classification of reduction types of curves}. \url{https://arxiv.org/abs/2512.09082}.
\item Literature assessment dated exactly \texttt{2026-08-22}: supplied status \texttt{partially\_solved}. The supplied assessment says that semistable reduction guarantees existence after a finite extension, that explicit algorithms cover important classes including hyperelliptic curves and degree-\(p\) covers, and that computation of the model and Galois action in the stated arbitrary-genus generality is incomplete. This is supplied dated curation derived from the OWR record, not an independently verified literature assessment; the cited claims, their hypotheses, and time-sensitive status require verification.
\item The excerpt refers to references [1], [2], and [16], without supplying their bibliographic entries. Reference [1] concerns joint work with Irene Bouw on \(L\)-factors and conductor exponents of superelliptic curves with a ``tame'' exponent. Reference [2] concerns joint work with Duc Do and Irene Bouw on Weil representations for curves with reduction of compact type. Reference [16] and Eliza Lorenzo Garcia's conference talk are cited for an alternative approach for smooth plane quartics. Tim Dokchitser's conference talk is mentioned concerning computation of regular models as an alternative for some applications. No additional citations or URLs are inferred.
\end{itemize}
\end{document}
